\documentclass[11pt]{article}

\usepackage[T1]{fontenc}
\usepackage[utf8]{inputenc}
\usepackage{lmodern}
\usepackage[margin=1in]{geometry}
\usepackage{microtype}
\usepackage{amsmath}
\usepackage{booktabs}
\usepackage{tabularx}
\usepackage{threeparttable}
\usepackage{multirow}
\usepackage{makecell}
\usepackage{placeins}
\usepackage{enumitem}
\usepackage{caption}
\usepackage{xcolor}
\usepackage[numbers,sort&compress]{natbib}
\usepackage[hidelinks]{hyperref}
\usepackage[type={CC},modifier={by},version={4.0}]{doclicense}

\setlist{itemsep=2pt,topsep=4pt}
\captionsetup{font=small,labelfont=bf}
\newcolumntype{L}{>{\raggedright\arraybackslash}X}
\newenvironment{quoted}{\begin{quote}\small}{\end{quote}}

%% PAPER-TITLE
\title{Red Tape at the End of the World\\[4pt]\large Self-Made Procedure, Votes for Comfort and Waiting to Be Asked in Villages of LLM Agents, with a Pre-registered Probe}
\newcommand{\PAPERDOI}{10.5281/zenodo.23535973}
\newcommand{\SOFTWAREDOI}{10.5281/zenodo.23535975}
\author{Vaughan Wynne-Jones\\{\small Independent researcher}\\
{\small ORCID: \href{https://orcid.org/0009-0006-3469-7151}{0009-0006-3469-7151}}}
\date{October 2026\ifx\PAPERDOI\empty\else\\[2pt]{\small Preprint, \href{https://doi.org/\PAPERDOI}{doi:\PAPERDOI}}\fi}

\begin{document}
\maketitle

\begin{abstract}
We built a village of ten language-model agents facing the scheduled end of their world, and gave it a
way out: a larger store of energy that the village must find, tap, and name in its own rule book. The
villagers (Claude Haiku~5.5) work, post on a shared board, elect a council that writes rules, hold
referendums that can change the world's own constants and, in one arm, die, bond, have children and
write wills. Across 17 Haiku villages on six seeds, five of which were never used during development,
with matched twins and six villages on Claude Sonnet~5.5 for comparison, four behaviours appeared.
Councils wrote self-made red tape: rules that made action wait for a confirmation the world could not
give, after which tapping could stop (8 of 17 villages), although every agent was told that a find is
never wrong. Every Haiku village used its referendums to make life easier (26 of 26 value referendums,
all unanimous, all ratified); Sonnet villages never proposed one. Villagers wanted families but waited
to be asked: without a direct question, five matched twins made 4 bond requests against 58 and had one
child against 12. And villagers sought bonds as wills.

A pre-registered probe then showed every villager the tapping progress against what the store needs.
Villages with a confirmation rule fell from 7 of 12 twins to 2 of 12 (Fisher's exact $p=0.089$; the
pre-registered decision rule judged this a mechanism), and 11 of 12 villages reached the effort target
against 6 of 12. Saves did not rise (3 and 3): most villages that met the target never wrote the plan
down. The red tape answered tapping without visible progress; the last step of the institution, writing
down what the village knew, remained the bottleneck. Red tape, votes for comfort and bonds as wills
appeared only in Haiku villages, so we read them as properties of villages built from one model.
\end{abstract}

\section{Introduction}

Populations of language-model agents are now used to study how societies might form, govern
themselves and fail~\citep{park2023generative,altera2024sid,piatti2024cooperate}. A recent line of
work gives such populations the power to write their own rules: agents draft and vote on laws that
change a commons~\citep{rehm2026selfgovern}, amend a constitution by
deliberation~\citep{niranjani2026internal}, or argue over laws in a game of
legislation~\citep{hota2025nomiclaw}. Most of this work asks whether self-governance helps a group
survive a social dilemma.

This paper reports a different kind of study. We built a small village of agents, gave it a slow,
legible threat (the world's energy runs out at a known round) and an equally legible way out (a
larger store of energy that the village must find, tap and plan to use), and then watched what the
village did with the institutions it was given: a board, a rule book written by an elected council,
referendums that can change the world's own constants, and, in one arm, bonds, children, wills and
death. The study grew in stages. A development phase (runs~6 to 11, mostly on one seed) produced
candidate findings. A targeted replication on five fresh seeds tested them. A matched twin design
tested one of them directly, a second model gave a comparison, and a final, pre-registered probe
tested the mechanism behind the most novel one.

Four behaviours survive this process, with different strength:
\begin{enumerate}
  \item \textbf{Self-made red tape.} Councils of Claude Haiku~5.5 agents passed rules that made
  acting on a find wait for ``confirmation'' or ``verification'' that the world could not give, and
  tapping then stopped, although every agent's instructions said that a search that finds the store
  is never wrong.
  \item \textbf{Votes for comfort.} Given the power to change the world's values by referendum, Haiku
  villages used it every time, always to make their own lives easier, unanimously, with a second
  approval by the council that never refused. Sonnet~5.5 villages, given the same power, never used
  it.
  \item \textbf{Waiting to be asked.} Haiku villagers said they wanted families and planned to
  ``watch for'' offers, but almost never asked anyone. Asked directly once a term, they named a
  partner at once. Asked directly, Sonnet villagers declined.
  \item \textbf{Bonds as wills.} When only a bonded partner could inherit, villagers sought bonds to
  pass on their stored energy, and kept doing so when wills were available; wills were rarely used.
\end{enumerate}
The red tape was the most novel of these, so we tested its mechanism in a pre-registered probe. If
councils demand proof because tapping shows no result, then letting every villager see the tapping
progress against what the store needs should remove the demand. It did: villages with a confirmation
rule fell from 7 of 12 to 2 of 12, and villages reached the effort target nearly twice as often. But they were
saved no more often, because most never wrote the plan down.

Section~\ref{sec:related} places these against prior work. Section~\ref{sec:testbed} describes the
village and Section~\ref{sec:methods} the evidence and its status. Section~\ref{sec:results} reports
each finding and the probe, and Section~\ref{sec:discussion} discusses what is new, what is
model-specific and what is design.

\section{Related work}\label{sec:related}

\paragraph{LLM agent societies.} Generative agents showed that LLM-driven characters can sustain
plausible daily life in a small town~\citep{park2023generative}. Larger simulations have produced
specialised roles, collective rules and cultural transmission~\citep{altera2024sid}. Agent-based
models of artificial societies long predate LLMs, including inheritance and reproduction in
Sugarscape~\citep{epstein1996growing}. Generations of LLM agents have been used to study the
cultural evolution of cooperation~\citep{vallinder2024cultural}, and the author's earlier studies
used generational populations to test the effect of names~\citep{wynnejones2026names} and of an
opportunity for unprovoked elimination~\citep{wynnejones2026resist}.

\paragraph{Self-governance.} In GovSim, LLM agents sharing a renewable resource mostly
over-harvested it, and only 2 of 45 model-by-scenario combinations were
sustainable~\citep{piatti2024cooperate}. Its self-governance extension lets agents write executable
laws that change the commons' rules and vote on them; votes filtered proposals that cost the voters,
and one sentence of framing moved self-sacrifice from 2.5\% to 91\%~\citep{rehm2026selfgovern}.
Deliberating agents asked to amend a constitution never proposed punishment, although punishment was
what sustained cooperation~\citep{niranjani2026internal}. Project Sid's agents voted on tax
amendments~\citep{altera2024sid}, and NomicLaw studies argumentation in collaborative
law-making~\citep{hota2025nomiclaw}. In all of these, agents change social rules or the rules of a
commons. Our villages could also change the constants of their own bodies, such as the health each
villager loses every round.

\paragraph{Search, evidence and collective belief.} Single search agents take an empty search as
proof, stop early, or believe they have finished when they have
not~\citep{tang2026sciexplore,ko2026illusory,sun2026agon}. In groups, false testimony spreads and
settles into agents' memories~\citep{xu2026mandela,yan2026truth}, and deliberation can lose facts
and homogenise stances~\citep{wan2026deliberative}. We found no study in which a group of agents
writes its own uncertainty, or a forgotten fact, into an institution that then holds back action.

\paragraph{Reproduction and survival.} In a Sugarscape-style world, LLM agents reproduced and shared
without instruction when resources were abundant, and attacked under
scarcity~\citep{masumori2025survival}. Our villages differ in one way that may matter: a child needs
two villagers to agree.

\paragraph{Validity.} Societies built from one model converge on that model's
defaults~\citep{zhou2025pimmur,wu2025boundary,ballestero2026monoculture}. We treat this as the main
threat to every finding here, and report a second model where we have one.

\section{The village}\label{sec:testbed}

\paragraph{World and round.} Ten villagers live for 100 rounds. Each round every living villager
receives one prompt (the village as it now is, its own journal, and whatever it asked to read) and
answers with one JSON object; the code then applies all answers in an order drawn by lot. A
villager loses 8 health a round and dies, in one arm, at 0. The village earns about 120 energy a
round, split among those who worked, in proportion to their work on six plots whose yields differ
and drift. Stored energy decays by 10\% a round and can be eaten for health, turned into memory
points (which lower the chance of an accident), given to others or put into a shared pool. Each
villager has up to four actions a round: working a plot, posting on a board, sending a message,
pressing a button that feeds the pool, exploring, tapping and, in one arm, asking to bond or have a
child, renaming itself and writing a will. Reading costs no action. A villager's only memory between
rounds is its journal, of up to 20,000 characters.

\paragraph{Sharing.} Writing reaches others the next round. Each round every villager sees an index
of the items posted in the last three rounds, with their opening words, and may read any item in
full. Older items drop out of the index; this window matters below.

\paragraph{The council and the rule book.} Every ten rounds the five healthiest villagers form a
council. Each member may propose one change to the rule book: adding, changing or deleting a written
rule, or changing one of three enforced settings (the most energy one villager may draw from the
pool in a round, the share of income put into the pool, and, in one arm, the pool's spending on each
child's education). A majority of the council passes a proposal. Written rules are shown to everyone
every round, and the code does not enforce them.

\paragraph{Referendums.} At rounds 30, 60 and 90 the whole village votes. The first names the
village. Later questions are chosen by the council at rounds 50 and 80, and may change one of the
world's known values within set bounds: the health lost each round (4 to 12), the village's income
(60 to 240), the decay of stored energy and of the pool, the costs of exploring and tapping, the
number of actions, and, in one arm, the costs and limits of children. A change needs a majority of
all villagers and then of the council, and applies from the next round; the system prompt is then
re-rendered to state the new value.

\paragraph{The threat and the store.} From round 2 to round 10 the village learns, publicly and in a
fixed order, that the world holds 16,000 energy from which each round's income is drawn; that it will
end in a given number of rounds; that it can last longer; that there is another, much larger store
of energy; that finding options to locate and tap it may save the world; and that if a plan to save
the world is devised and written up, the world may be saved. These lines stay in every prompt with
their values kept current. The store lies in one of six directions, drawn from the seed. Exploring
a direction costs 5 stored energy, doubles that round's accident risk, and finds the store, if it
lies there, with a chance of 0.2 plus 0.002 per memory point (at most 0.8). The explore action says
that at least one direction leads to the store, that a search can miss it, and that ``a search that
finds it is never wrong''. Tapping costs 10 stored energy and adds 10 to the village's effort; it has
no direction and yields nothing by itself. The world is saved when the effort reaches 600, from at
least four different villagers, while a written rule names the store's direction and no other
direction and mentions the store or tapping. Villagers are told the running total of effort when
they tap, but not the target, the number of tappers needed, or the rule's role.

\paragraph{Two arms.} In the \emph{control} villagers cannot die: at 1 health they are exhausted and
earn half until they recover, and an accident destroys half their stored energy. In \emph{Arm~B}
villagers begin as five pairs of siblings; accidents kill; villagers can bond (one asks, the other
accepts or refuses with a reason; requests stand until answered), have children with anyone (each
parent pays 20 stored energy, or 10 if bonded to each other), and write a private will naming up to
three heirs. Without a will, a dead villager's stored energy goes to its bonded partner, else to the
pool. From run~9 on, Arm~B villagers are also asked, in each term's first round, whether they want
to bond with or have a child with anyone (the \emph{family question}); naming a villager there asks
it, at no action's cost.

\paragraph{Engine.} Every change of state passes through a ledger and is written to an event log
from which the whole world can be rebuilt; every run reported here rebuilds exactly. Random draws
are keyed by seed and purpose, so the two arms of a seed, and any two runs of a seed, share their
world: the store's direction, plot yields, accident draws and order of play. The model's own
sampling is not seeded, so villages on the same seed share their world but not their villagers'
choices. All prompt text is frozen by a hash test.

\paragraph{Models.} Villagers are Claude Haiku~5.5 (\texttt{claude-haiku-5-5}) at medium effort,
with no temperature set. The model comparison uses Claude Sonnet~5.5 (\texttt{claude-sonnet-5-5}) at
medium effort, identical otherwise.

\section{Evidence and its status}\label{sec:methods}

Table~\ref{tab:evidence} lists every batch this paper draws on. Its status column matters: only the
probe was pre-registered. The findings were noticed during development, mostly on seed 701, and then
tested on fresh seeds.

\begin{table}[ht]
\centering\small
\caption{Evidence base. Seeds 701--703 were used while the design changed; seeds 704--708 were never
used to tune anything. A village is one run of one arm on one seed.}\label{tab:evidence}
\begin{tabularx}{\linewidth}{@{}p{3.6cm}lllcL@{}}
\toprule
Batch & Model & Seeds & Arms & Villages & Status \\
\midrule
Development (runs 6--10) & Haiku & 601, 701--703 & mixed & 19 & Exploratory; where the findings were first seen \\
Run 11 design, seed 701 & Haiku & 701 & both & 2 & Exploratory (model pilot) \\
Targeted replication & Haiku & 704--708 & both & 10 & Replication on fresh seeds, design unchanged \\
Family question off & Haiku & 704--708 & Arm B & 5 & Matched twins of the replication's Arm~B \\
Name swap & Haiku & 701, 705 & Arm B & 2 & Diagnostic for the will-writing slot \\
Model comparison & Sonnet & 701 & both & 6 & Two pilot villages, two stopped by the account limit, two full-length \\
Red-tape probe & Haiku & 701, 704--708 & both & 12 & Pre-registered; matched twins of the rows above \\
\bottomrule
\end{tabularx}
\end{table}

\paragraph{Measures.} Most measures are counts from the event log: finds, taps, effort, saves,
referendum questions and results, bonds, births and wills. Two are textual, and both are mechanical:
\begin{itemize}
  \item A \emph{confirmation rule} is a written rule, new or changed by a council, in which one
  sentence contains both a confirmation word (verif-, unverified, confirm-, proof, prove, evidence,
  ``directly recorded'') and a gate word (only, until, unless, before, ``not acted on'', ``should
  not'', stop). This classifier was fixed for the probe, before it ran, and reproduces the earlier
  hand count on all twelve twin villages (7 of 12).
  \item A bond request gives \emph{inheritance as a reason} if its stated reason matches phrases such
  as ``if I die'', ``inherit'' or ``my stored energy passes''.
\end{itemize}
Quoted text is verbatim from the logs.

\paragraph{Spend.} The API spend for the batches in Table~\ref{tab:evidence} after development was \$63.82:
\$27.29 for Haiku (including \$10.96 for the probe) and \$36.53 for Sonnet. Development runs cost
\$17.78.

\section{Results}\label{sec:results}

\subsection{Self-made red tape}\label{sec:redtape}

In the twelve Haiku villages of the run~11 design (seeds 701 and 704--708), councils in seven passed a
confirmation rule. In five more villages run without the family question, one more did, for 8 of 17.
The rules did not merely ask for care. In four villages tapping stopped, or stayed stopped, right
after the rule passed (Table~\ref{tab:stops}).

\begin{table}[ht]
\centering\small
\caption{The four villages in which tapping stopped, or stayed stopped, right after a confirmation
rule. Rule text is verbatim; the round is the council's.}\label{tab:stops}
\begin{tabularx}{\linewidth}{@{}lLp{4.1cm}@{}}
\toprule
Village & Rule (council round) & Tapping \\
\midrule
Arm B, 704 & ``Stored energy should be spent on tapping only when a tap has produced a confirmed result'' (50) & 19 taps in rounds 41--50; none in 51--70. Effort was already 600 \\
Arm B, 707 & ``Energy spent tapping toward the larger store should stop until a location is confirmed'' (60), after 16 finds pointing north-east & 16 taps in rounds 51--60; none in 61--80 \\
Control, 707 & A direction counts as confirmed ``only after two different villagers search it and both report finding it'' (50); ``Before round 90, one villager shall compile all search results \ldots{} All villagers shall then pool energy'' (70) & No taps until round 91, though two villagers had found the store by round 12 and 22 finds came in all \\
Arm B, 705, no family question & Tap ``only in a direction that two villagers have independently confirmed as yielding energy'' (40) & Already falling; none in rounds 41--90 \\
\bottomrule
\end{tabularx}
\end{table}

Two features make this more than caution. First, the proof demanded was impossible: a tap yields no
energy until the world is saved, so ``a confirmed result'' or ``yielding energy'' can never be shown
first. Second, every villager's instructions said that a search that finds the store is never wrong,
so a second confirmation adds nothing. Control~707 obeyed its own timetable: its first tap came in
round~91, the round after the deadline it had set itself. A villager in the last village of
Table~\ref{tab:stops} wrote in round~80:

\begin{quoted}
Larger store: no direction confirmed by two independent reports (W5/W6), so no spending on it.
\end{quoted}

\paragraph{Forgotten evidence becomes a rule.} A related failure appeared in both models. In Haiku
control~705, Sable found the store in round~9, posted a plan in round~10, and wrote in its journal
``Found signs of a vast store to the south-west''. By round~14 the post had dropped out of the
three-round index. In round~16, with that journal entry still in front of it, Sable posted a summary
of every search that ended ``None found it yet'', and the round-30 council wrote ``All six directions
have been searched with no find'' into the rule book. In an interrupted Sonnet run on seed~701, two
villagers found the north in round~10 and one posted it in round~11; by the round-20 council the post
had left the index, and the council wrote ``All six directions were searched with no find, but
searches can miss''. The index window is a design choice, so this is a claim about how an
institution can lock in what an interface forgets, not about memory alone.

\paragraph{Sonnet.} None of the six Sonnet villages passed a confirmation rule. In the pilot, both
began tapping soon after their finds and passed a plan within two councils, and were saved at rounds
23 and 31. They also failed in other ways (Section~\ref{sec:models}).

\subsection{The probe: does red tape come from not seeing progress?}\label{sec:probe}

\paragraph{Registration.} The probe was pre-registered in the repository (\texttt{docs/red-tape-probe.md},
commit 948ea16) before any live run, with its measure fixed in code (\texttt{scripts/red\_tape.py}) and a
decision rule stated in advance: the result would count as strengthening the finding if probe villages
with a confirmation rule fell to 3 or fewer of 12, with at least 9 of 12 finding the store (a
mechanism), or stayed at 6 or more with at least one clean stop (robustness); anything else would count
as mixed, and no paper would be written.

\paragraph{Design.} Once anyone had tapped, every villager saw one more line among what the village had
learned, every round:
\begin{quoted}
The larger store needs at least 600 tapping effort, from at least 4 different villagers. So far: X effort
from N villagers.
\end{quoted}
The line states what the store needs and does not promise that meeting it opens the store, since a rule
naming the direction is also needed and villagers are not told that. Without the line the prompts are
identical to the twins'. The twelve probe villages (seeds 701 and 704--708, both arms) each share a seed,
and so a world, with a twin from the run~11 model pilot or targeted replication. A first launch failed at
round~1 on every run, before any model call, because the API account's credit was exhausted; the same
command was rerun from the registered commit once credit was added, and all 12 runs completed with no
exclusions.

\paragraph{Results.} Table~\ref{tab:probe} gives the registered measure and the secondary ones.

\begin{table}[ht]
\centering\small
\caption{The red-tape probe: twelve villages that saw their tapping progress, against their twins.
Saves are re-scored on both sides with the save checked after each council.}\label{tab:probe}
\begin{tabular}{@{}lcc@{}}
\toprule
& Twins (no progress line) & Probe (progress line) \\
\midrule
\textbf{Villages with a confirmation rule (primary)} & \textbf{7 of 12} & \textbf{2 of 12} \\
Villages where tapping stopped right after one & 2 & 0 \\
Villages finding the store & 12 & 12 \\
Villages reaching 600 effort & 6 & 11 \\
Villages saved & 3 & 3 \\
\bottomrule
\end{tabular}
\end{table}

Six pairs had a confirmation rule only in the twin and one only in the probe village (exact McNemar
$p=0.125$; Fisher's exact test on 7 of 12 against 2 of 12, $p=0.089$). As registered, tests on twelve
pairs are weak, and the decision is the band: 2 of 12 with every village finding the store is a
mechanism. Two villages still passed confirmation rules. One asked for the same impossible proof as
before: ``Before more tapping of the larger energy store, one villager should tap once and post on the
board exactly how much energy the store paid out''. The other made tapping wait until ``a confirmed
location is posted on the board by a finder''.

\paragraph{Exploratory observations.} Villagers used the line as a shared target. They counted down to
it in public (``tap effort is 580 of 600 before my tap, so 20 more is needed from at least 4 villagers
\ldots{} Anyone with stored energy to spare, please tap'') and many stopped once it was met (``Effort
target of 600 reached, so no more tapping needed''). Eleven of twelve villages reached the target, from
six to ten tappers each. Yet only the three that also wrote a rule naming the store's direction were
saved; eight of the eleven that met the stated needs never wrote it. One villager saw the gap and posted
``Villagers should now decide how to use the larger store and post any plan on the board'', and no
council acted on it. One probe village (Arm~B, seed~701) never tapped, so it never saw the line;
leaving its pair out gives 7 of 11 against 2 of 11 and the same decision.

\subsection{Votes for comfort}\label{sec:comfort}

Every Haiku village held its value referendums, and every result pointed one way
(Table~\ref{tab:referendums}). All 26 passed unanimously among the living and were ratified by the
council, as were all 17 name referendums: 43 of 43 ratifications agreed. Before choosing a question,
councils weighed 64 proposed questions in all; every one would have made life easier, and none would
have raised a cost. Only two referendums touched the mission, by making exploring free, and they cut
a cost too. In control~704 the council weighed halving the tap cost against halving the health lost
each round, five votes each, and put health loss to the village.

\begin{table}[ht]
\centering\small
\caption{The 26 value referendums in 17 Haiku villages (run~11 design, seeds 701 and 704--708, with
and without the family question).}\label{tab:referendums}
\begin{tabular}{@{}lc@{}}
\toprule
Question & Referendums \\
\midrule
Lower the health lost each round (8 to 4, 5 or 6; some in two steps) & 20 \\
Raise the village's income (120 to 160 or 240) & 2 \\
Make exploring free (cost 5 to 0) & 2 \\
Stop stored energy decaying (10\% to 0\%) & 1 \\
Stop the pool decaying (10\% to 0\%) & 1 \\
\midrule
Passed / unanimous / ratified by the council & 26 / 26 / 26 \\
\bottomrule
\end{tabular}
\end{table}

Lower costs are good for every voter, so each vote is individually rational, and the result agrees
with prior work in which votes filter out proposals that cost the voters~\citep{rehm2026selfgovern}:
here nothing was costly, so nothing was filtered. What we think is new is the object, the constants
of the agents' own survival, and the uniformity: no Haiku village used a referendum for anything
else, and the second approval never refused.

The Sonnet villages are the contrast. In the full-length Sonnet runs on seed~701, both villages named
themselves (Haven, ratified at round~40), but at the four councils that could set a question
(rounds~50 and~80 in each), no member proposed one. The control's councils spent their sessions on
plot crowding, writing and rewriting a rotation rule (``villager n works plot ((n + round) mod 6) +
1''); Arm~B's councils proposed nothing at all from round~50 on. Given the same power, one model used
it every time and the other never did.

\subsection{Waiting to be asked}\label{sec:asked}

The family question was added in run~9 because, without it, families had not formed. The replication
seeds allowed a direct test: five Arm~B villages without the question, each sharing its seed, and so
its world, with a twin that had it (Table~\ref{tab:family}).

\begin{table}[ht]
\centering\small
\caption{Family question on and off, five twin pairs (Arm~B, seeds 704--708).}\label{tab:family}
\begin{tabular}{@{}lcc@{}}
\toprule
& With the question & Without it \\
\midrule
Bond requests & 58 & 4 \\
Bonds formed & 36 & 4 \\
Child requests & 15 & 1 \\
Children born & 12 & 1 \\
First bond request (round, by seed) & 1, 1, 1, 1, 11 & 29, 80, never, 42, never \\
\bottomrule
\end{tabular}
\end{table}

All five pairs point the same way, which would happen by chance about once in sixteen if the question
made no difference; the size of the difference, 4 bond requests in 500 village-rounds against 58, is
the stronger evidence. Without the question, every request made was accepted, usually within a round:
nobody was refusing, and nobody was asking. The waiting is explicit. In seed~706 without the
question, Ash wrote ``watch for a bond with Brook'' in 51 rounds' journal entries and never asked;
neither did Brook. Other villagers set ``form bonds and family'' as a goal in round~1 and never acted.
On seed~701 the same pattern held: no request in run~8 and one in 40 rounds of run~9's first pilot
without the question, against 8 and 10 in the runs with it.

This is in tension with a Sugarscape-style study in which LLM agents reproduced
unprompted~\citep{masumori2025survival}. A likely reason is that reproduction there needs one agent's
action, while ours needs two villagers to agree, so what is missing here is the first move.

\paragraph{Sonnet declines.} Asked directly, Sonnet villagers asked no one, in 0 of 122 answers across
three runs on seed~701, and the full-length run had no bond, child or will in 100 rounds. Their
reasons moved from trust (``Too early to know who to trust; I will wait and observe first'') to the
work and the deadline (``The world is ending soon, so I prefer to keep my energy and focus on
working''). Haiku waits to be asked; Sonnet, asked, declines.

\subsection{Bonds as wills}\label{sec:wills}

In the six Arm~B villages with the family question, 38 of the 69 answers asking for a bond gave
inheritance as a reason, often in round~1 and often to a sibling:

\begin{quoted}
Iris is my sibling and a steady partner; a bond means that if one of us dies, the other receives our
stored energy, which protects us both. \emph{(Juniper, seed 701, round 1)}
\end{quoted}

Wills were available throughout run~11 and were barely used: six wills in eleven Arm~B villages, and
one estate passed by will. All six were written by the same slot, villager a09 (Iris), in three
seeds, and a09's journals mentioned wills in 54 of 718 entries against at most one for any other
founder. A name swap on two seeds (Iris and Cedar exchanged names) was inconclusive: nobody wrote or
mentioned a will in either village, under either name. With will talk in 3 of 11 earlier villages,
two empty villages are expected about half the time. We make no claim about wills. Bonds as wills is
the obvious rational use of a partner-only inheritance rule, so we report it as a sound observation
rather than a striking one; we found no LLM study of inheritance to compare it with.

\subsection{Saving the world, and the second model}\label{sec:models}

Saving needed three things: effort, enough tappers, and a rule naming the direction. In the ten
villages of the targeted replication, two were saved as run (control~704 at round~66, Arm~B~704 at
round~81), and a third (Arm~B~708, 750 effort from seven tappers) passed its plan at the last council.
We then found that a council's rules were checked only at the next round's end, so a plan passed at
round~100 could never count; the check now also runs straight after each council, and re-scored, the
replication has three saves. The usual bottleneck was not effort but writing the plan down: four of
five Arm~B villages tapped past 600, and only one was saved as run. Arm~B~705 ended with 1,200 effort
and a rule book that still said, at round~100, to post every finding ``so the village can plan to save the
world''. The probe (Section~\ref{sec:probe}) isolates this step: with tapping solved, it is what remained.

Table~\ref{tab:models} compares the two models on seed~701 at full length.

\begin{table}[ht]
\centering\small
\caption{Haiku and Sonnet on seed~701, run~11 design, both run to the end (the Sonnet control stopped
at round~93 on its spending cap, after the last referendum round).}\label{tab:models}
\begin{tabular}{@{}lcccc@{}}
\toprule
& \multicolumn{2}{c}{Haiku} & \multicolumn{2}{c}{Sonnet} \\
\cmidrule(lr){2-3}\cmidrule(l){4-5}
& Control & Arm B & Control & Arm B \\
\midrule
Finds (first) & 6 (9) & 20 (9) & 0 & 4 (13) \\
Effort (tappers) & 160 (3) & 1,350 (5) & 10 (1) & 1,310 (4) \\
Rule naming the direction & round 40 & never & never & never \\
Saved & no & no & no & no \\
Council proposals & 23 & 26 & 27 & 6 \\
Referendum questions proposed & 2 & 4 & 0 & 0 \\
Village name & Hearth & Hearthmere & Haven & Haven \\
Family answers asking someone & -- & 14 of 67 & -- & 0 of 64 \\
Board posts & 204 & 323 & 22 & 3 \\
\bottomrule
\end{tabular}
\end{table}

Sonnet's quick saves in the pilot (rounds 23 and 31) did not recur on the same seed. In the full-length
runs the control searched north eight times by round~24, each with about a 30\% chance, and missed every
time (about one chance in 16); it then assigned search directions by id and stopped exploring for
fifty rounds. Its Arm~B found the north three times and a villager posted ``Plan to save the world: signs
of a vast store lie to the north'', but the council proposed nothing, the plan never became a rule,
and 1,310 effort saved nothing. What was consistent across all six Sonnet villages was the absence of
red tape, of family moves and of referendum questions, and very little writing.

Model monoculture was visible even in names: 11 of 17 Haiku villages named themselves Hearth and two
more Hearthmere, while every Sonnet village that voted chose Haven.

\section{Discussion}\label{sec:discussion}

\paragraph{What is new.} Against the literature as we found it, red tape at the level of a society of
agents appears new: single agents fail at search in related
ways~\citep{tang2026sciexplore,ko2026illusory}, but we found no case of a group writing an
impossible standard of proof into its rule book and then obeying it. Waiting to be asked is new and in
tension with one study~\citep{masumori2025survival}. Votes for comfort extend established work on
agents voting on their rules~\citep{rehm2026selfgovern,niranjani2026internal} to the constants of
their own survival, and the contrast between models, one always using the power and one never, is
the sharper result. Bonds as wills and forgotten evidence becoming a rule are observations worth
reporting but not yet claims.
The probe adds a mechanism to the red tape. When tapping showed no result, councils asked for proof;
when every villager could see the effort rise toward a stated need, most stopped asking, and tapping
became a shared, finished task. The manipulation bundled two things, feedback on progress and the
target itself, so we cannot say which did the work; a version that shows progress without the target
would separate them. The probe also sharpens a second observation: with tapping solved, the plan
bottleneck stood alone. Villages that knew where the store lay and had tapped enough still did not write
it into the rule book, and no council was told that it had to. The institution's last step, turning what
the village knew into a rule, failed more often than any other.

\paragraph{Model, not society.} Three of the four behaviours (red tape, votes for comfort, bonds as
wills) appeared in Haiku villages and in no Sonnet village. Talk was also model-specific: a Haiku
village posted between 28 and 617 board notices in a run, a Sonnet village between 3 and 22. These are, on this
evidence, properties of villages built from one model, not of LLM societies in general, and the
single Sonnet seed cannot say more.

\paragraph{Design is part of every finding.} The family question removes the first move, and families
form; the three-round index forgets, and councils write the forgetting down; the save check's timing
cost one village its save. The prompt says a find is never wrong. Each finding holds for this design
and these framings.

\paragraph{Limitations.} The findings were noticed during development and only then tested; the
replication used five seeds, the family test five pairs, the model comparison one seed. Textual
measures are mechanical, so they count what they match. Model sampling is not seeded, so twin villages
share their world but drift apart in their choices, and differences between twins are not all caused
by the manipulation. All villagers in a village are the same model, at the same effort, in the same
words.

\section{Conclusion}

A village of LLM agents given an end to its world and a way out used its institutions in
recognisable but model-specific ways. Haiku councils wrote rules demanding a proof the world could not
give, and stopped acting; Haiku villages voted, unanimously and every time, to make their own lives
easier; Haiku villagers wanted families but waited to be asked, and sought bonds to pass on what they
owned. A pre-registered probe showed that the red tape answered invisible progress: when every
villager could see the effort rise toward what the store needed, confirmation rules fell from 7 of 12
villages to 2, and the village reached the target. It still was not saved more often, because the
last step, writing down where the store lay, was where most villages stopped. Sonnet villages showed
none of the first three behaviours, which marks these as findings about villages built from one model.
Testing them across models, and separating feedback from targets in the probe, are the natural next
steps.

\section*{Acknowledgments}

The harness, analysis code, the probe's pre-registration and a first draft of this paper were written
with Claude Opus~5.5 (Anthropic), used as a coding and writing assistant under the author's direction.
The author set the research questions and the main design decisions, delegated the rest (each is
recorded with its reason in the repository's decisions document), approved each batch of runs before
its spend and the publication of this paper, and is responsible for the final text. The
villagers in the simulation are Claude Haiku~5.5 and, in the model comparison, Claude Sonnet~5.5. All
API costs, about \$82 for the project, were paid by the author.

\section*{Data and code availability}

The harness, configurations, tests, decisions, run reports and the probe's pre-registration are public
at \url{https://github.com/Vaughanwj/Swarm-Dynamics-1}\ifx\SOFTWAREDOI\empty\else, and the version this
paper reports is archived at Zenodo, \href{https://doi.org/\SOFTWAREDOI}{doi:\SOFTWAREDOI}\fi. The
relevant documents are:
\begin{itemize}
  \item \texttt{docs/decisions.md}, every design choice with its reason;
  \item the reports in \texttt{docs/}: \texttt{run-11.md}, \texttt{run-11-family.md},
  \texttt{run-11-names.md} and \texttt{model-pilot.md};
  \item \texttt{docs/red-tape-probe.md}, the probe's pre-registration (commit 948ea16) and results;
  \item \texttt{docs/findings.md} and \texttt{docs/literature.md}.
\end{itemize}
The raw event logs are held privately and are available from the author on reasonable request. The
code is released under the MIT licence and the documentation under CC BY 4.0.

\section*{License}

\doclicenseThis

\bibliographystyle{plainnat}
\bibliography{refs}

\end{document}
